\section{Conclusion}
\label{sec:conclusion}

We presented \gore{}, a minimal five-primitive formal language that makes the
search tree a first-class syntactic object, together with a three-stage training
pipeline that leverages the structural isomorphism between \gore{} and
Lean\,/\,Coq tactic combinators.
The central empirical finding is that dense step-level reward from the \gore{}
interpreter (acting as a Process Reward Model) substantially improves backtrack
accuracy over sparse final-answer reward, confirming that the locus of the
learning signal matters for search-requiring tasks.
The Spectral-R1 analysis provides mechanistic evidence that the \textsc{Fork}
and \textsc{Cut} primitives are encoded in a distinguishable representational
subspace, suggesting that backtracking has a learnable internal structure.

\paragraph{Limitations.}
\gore{} is a synthetic environment; the gap between synthetic backtracking
performance and natural-language reasoning performance remains open.
The current interpreter reward does not account for semantically equivalent but
syntactically distinct traces (e.g., permuted \textsc{Fork} branches that yield
the same result), which may over-penalise valid alternative solutions.
Experiments are conducted at GPT-2 scale; behaviour at larger scales is unknown.

\paragraph{Future work.}
\begin{itemize}
  \item \textbf{YOCO integration.} For multi-turn \gore{} (recursive task
    decomposition), the YOCO decoder-decoder architecture \citep{sun2024yoco}
    is a natural fit: the self-decoder caches the static \gore{} program, and
    the cross-decoder generates the execution trace.
    This separation would allow the Spectral-R1 analysis to isolate program
    representations from trace representations.
  \item \textbf{Scaling.} Repeat E1--E3 at 1B and 7B parameter scales to
    determine whether the backtracking subspace becomes more or less separable
    with model capacity.
  \item \textbf{Constraint satisfaction.} Extend \gore{} with typed \callstmt{}
    and a constraint store to model constraint-satisfaction problems, bridging
    the gap to real-world planning and scheduling tasks.
  \item \textbf{Multi-agent \gore{}.} Treat each \callstmt{} as a sub-agent
    invocation, turning \gore{} into an orchestration language for multi-agent
    systems with verifiable dense reward at every node.
\end{itemize}

\todo{Add final acknowledgements and funding statement before camera-ready.}
